\documentclass[12pt,a4paper]{article}
\usepackage[spanish]{babel}
\usepackage[utf8]{inputenc}
\usepackage{amsmath}
\usepackage{amsfonts}
\usepackage{amssymb}
\usepackage{geometry}
\geometry{margin=2.5cm}

\title{\LARGE Trabajo Autónomo - Ayudantia\\ Elipse}
\author{}
\date{}

\begin{document}

\maketitle

\vspace{0.3cm}
\centering
\textbf{Preparado por:} \textit{Diego Solis S1, Joaquin Vasquez S2, Benjamin Romo S3}
\par
\vspace{0.7cm}

El siguiente listado corresponde a ejercicios referidos a \textbf{ELIPSE} del texto \textbf{[Lehmann]GeometriaAnalitica}

\begin{enumerate}
\item La ecuación de una elipse es \( x^2 + 4y^2 + 2x - 12y + 6 = 0 \).
Reducir esta ecuación a la forma ordinaria y determinar las coordenadas del centro,
de los vértices y de los focos; calcular las longitudes del eje mayor, del eje menor,
de cada lado recto y la excentricidad.
\item Reducir la ecuación dada a la segunda forma ordinaria de la ecuación de una elipse,
y determínense las coordenadas del centro, vértices y focos, las longitudes de los ejes
mayor y menor, y la de cada lado recto y la excentricidad:
\[ x^2 + 4y^2 - 6x + 16y + 21 = 0. \]
\item \[ 4x^2 + 9y^2 + 32x - 18y + 37 = 0. \]
\item \[ x^2 + 4y^2 - 10x - 40y + 109 = 0. \]
\item \[ 9x^2 + 4y^2 - 8y - 32 = 0. \]
\item Hallar la ecuación de la elipse cuyos vértices son los puntos \((4,0)\), \((-4,0)\),
y cuyos focos son los puntos \((3,0)\), \((-3,0)\).
\item Los vértices de una elipse son los puntos \((0,6)\), \((0,-6)\), y sus focos son
los puntos \((0,4)\), \((0,-4)\). Hallar su ecuación.
\item Hallar la ecuación de la elipse cuyos focos son los puntos \((2,0)\), \((-2,0)\),
y su excentricidad es igual a \( \frac{2}{3} \).
\item Los focos de una elipse son los puntos \((3,0)\), \((-3,0)\), y la longitud de uno
cualquiera de sus lados rectos es igual a 9. Hallar la ecuación de la elipse.
\item Hallar la ecuación y la excentricidad de la elipse que tiene su centro en el origen,
uno de sus vértices en el punto \((0,-7)\) y pasa por el punto \( \left( \sqrt{5},\ \frac{14}{3} \right) \).
\item Una elipse tiene su centro en el origen y su eje mayor coincide con el eje \(X\).
Hallar su ecuación sabiendo que pasa por los puntos \( (\sqrt{6},\ -1) \) y \( (2,\ \sqrt{2}) \).
\item Focos \((3,8)\) y \((3,2)\); longitud del eje mayor \(= 10\).
\item Vértices \((-3,-1)\) y \((5,-1)\); excentricidad \(= \frac{3}{4}\).
\item Vértices \((2,6)\) y \((2,-2)\); longitud del lado recto \(= 2\).
\item La ecuación de una familia de elipses es \( kx^2 + 4y^2 + 6x - 8y - 5 = 0 \).
Hallar las ecuaciones de aquellos elementos de la familia que tienen una excentricidad
igual a \( \frac{1}{2} \).
\item Hallar las longitudes de los radios vectores del punto \((2,1)\) de la elipse
\[ 9x^2 + y^2 - 18x - 2y + 1 = 0. \]
\item El punto medio de una cuerda de la elipse
\[ x^2 + 4y^2 - 6x - 8y - 3 = 0 \]
es el punto \((5,2)\). Hallar la ecuación de la cuerda.
\end{enumerate}

\end{document}